
\documentclass[10pt,a4paper]{article}

\usepackage[a4paper,margin=0.45in]{geometry}
\usepackage{enumitem}
\usepackage{array}
\usepackage{xcolor}
\usepackage{fancyhdr}
\usepackage{multicol}
\usepackage{amsmath}
\usepackage[T1]{fontenc}
\usepackage{lmodern}

\setlength{\parindent}{0pt}
\setlength{\parskip}{2pt}

\definecolor{univblue}{RGB}{20,60,110}

% -----------------------------------------------------
% PAGE STYLE
% -----------------------------------------------------
\pagestyle{empty}

% -----------------------------------------------------
% CUSTOM COMMANDS
% -----------------------------------------------------

\newcommand{\examheader}[1]{
\begin{center}
{\large\bfseries GRAPHIC ERA DEEMED TO BE UNIVERSITY}\\[-1mm]
{\small Department of Computer Science and Engineering}\\[1mm]
{\bfseries MID-SEMESTER PRACTICAL EXAMINATION}\\[-1mm]
{\small Operating Systems Laboratory}\\[1mm]
{\large\bfseries SET -- #1}
\end{center}

\vspace{-1mm}

\begin{tabular}{@{}p{0.33\textwidth}p{0.32\textwidth}p{0.31\textwidth}@{}}
\textbf{Name:} \hrulefill &
\textbf{Roll No.:} \hrulefill &
\textbf{Section:} \hrulefill
\end{tabular}

\vspace{1mm}
\hrule
\vspace{1mm}
}

\newcommand{\problem}[2]{
\textbf{Problem #1.} #2
}

% =====================================================
% DOCUMENT
% =====================================================
\begin{document}

% =====================================================
% SET 1
% =====================================================

\examheader{1}

\textbf{Programming Problems}

\begin{enumerate}[leftmargin=5mm,itemsep=2pt,label=\textbf{P\arabic*.}]

\item
Write a C program using \texttt{fork()} in which the
\textbf{parent process calculates the sum of odd numbers} and the
\textbf{child process calculates the sum of even numbers} from a
user-provided integer array. Display the PID and PPID of the processes.

\hfill\textbf{[6 Marks]}

\item
Write a C program to establish \textbf{Inter-Process Communication
using an unnamed pipe}. The child process should send a string to the
parent process through the pipe, and the parent should receive and
display the string.

\hfill\textbf{[6 Marks]}

\end{enumerate}

\vspace{1mm}

\textbf{Short Answer Questions \hfill (2 Marks Each)}

\begin{enumerate}[leftmargin=7mm,itemsep=1pt,label=\textbf{Q\arabic*.}]

\item What values are returned by \texttt{fork()} to the parent and
child processes?

\item What is the purpose of the \texttt{wait()} system call?

\item What is the difference between a \textbf{zombie process} and an
\textbf{orphan process}?

\item In \texttt{pipe(fd)}, explain the purpose of
\texttt{fd[0]} and \texttt{fd[1]}.

\end{enumerate}

\vspace{1mm}

\begin{center}
\textbf{Total Marks: 20}
\end{center}

% =====================================================
% DIVIDER
% =====================================================

\vspace{1mm}

\noindent
\makebox[\textwidth]{%
\dotfill\quad\textbf{\small CUT / SEPARATE HERE}\quad\dotfill
}

\vspace{1mm}

% =====================================================
% SET 2
% =====================================================

\examheader{2}

\textbf{Programming Problems}

\begin{enumerate}[leftmargin=5mm,itemsep=2pt,label=\textbf{P\arabic*.}]

\item
Write a C program using \texttt{fork()} and \texttt{wait()} where the
\textbf{child process prints numbers from 1 to 10} and the
\textbf{parent process prints numbers from 11 to 20}. Ensure that the
parent starts printing only after the child has completed execution.

\hfill\textbf{[6 Marks]}

\item
Write a C program using an \textbf{unnamed pipe} where the parent
process sends an integer to the child process. The child should receive
the integer, calculate its \textbf{square}, and display the result.

\hfill\textbf{[6 Marks]}

\end{enumerate}

\vspace{1mm}

\textbf{Short Answer Questions \hfill (2 Marks Each)}

\begin{enumerate}[leftmargin=7mm,itemsep=1pt,label=\textbf{Q\arabic*.}]

\item What is the difference between \texttt{getpid()} and
\texttt{getppid()}?

\item What happens to the parent process when it calls \texttt{wait()}
while its child is still executing?

\item What is an \textbf{unnamed pipe}, and why is it generally created
before calling \texttt{fork()}?

\item What happens when a process tries to read from an empty pipe while
the write end is still open?

\end{enumerate}

\vspace{1mm}

\begin{center}
\textbf{Total Marks: 20}
\end{center}

\end{document}